Los experimentos se ejecutaron en un equipo con Windows utilizando WSL2
(Ubuntu), \texttt{g++} con las opciones \texttt{-std=c++17 -O2 -Wall -Wextra}, Python 3,
NumPy y Matplotlib. Los programas midieron tiempo en microsegundos y memoria
residente en KB mediante \texttt{getrusage}. Cada algoritmo recibió una copia del
mismo input y se verificó el resultado antes de guardar la medición.

Para sorting se utilizaron los tamaños $n\in\{10,10^3,10^5\}$ con los tres
tipos de orden y dominios $D1,D7$, tres muestras por combinación, más una
entrada representativa $10^7_{\text{aleatorio},D7,a}$. En total se procesaron
55 entradas. Para matrices se utilizaron $n\in\{2^4,2^6,2^8,2^{10}\}$,
los tipos dispersa, diagonal y densa, dominios $D0,D10$ y tres muestras,
correspondientes a 72 pares de matrices. Los datos se generan con los
scripts de \texttt{code/sorting/scripts/} y
\texttt{code/matrix\_multiplication/scripts/}.
